% Generated draft: requires mathematical and source review.
\documentclass[11pt]{article}
\usepackage[margin=25mm]{geometry}
\usepackage{fontspec}
\setmainfont{Latin Modern Roman}
\usepackage{amsmath,amssymb,mathtools,unicode-math}
\setmathfont{Latin Modern Math}
\usepackage{enumitem,needspace,xurl,hyperref}
\setlength{\parindent}{0pt}
\setlength{\parskip}{5pt}
\setlength{\emergencystretch}{3em}
\newcommand{\N}{\mathbb{N}}
\newcommand{\Z}{\mathbb{Z}}
\newcommand{\Q}{\mathbb{Q}}
\newcommand{\R}{\mathbb{R}}
\newcommand{\C}{\mathbb{C}}
\newcommand{\F}{\mathbb{F}}
\newcommand{\A}{\mathbb{A}}
\newcommand{\PP}{\mathbb{P}}
\newcommand{\E}{\mathbb{E}}
\newcommand{\eps}{\varepsilon}
\newcommand{\dd}{\,\mathrm{d}}
\begin{document}
\section*{UnsolvedMath ID: 30000320}
\textbf{Source problem code:} catalogue entry
\subsection{Definitions and mathematical statement}
\textbf{Rhetorical transition, not a verified standalone open problem.} The original extraction says to determine the structure of the primitive-central cohomology algebra $\mathrm{PC}\,H^*(G)$ of a finite group $G$. Its verification field marks this unrecoverable: the official abstract immediately answers the rhetorical question with structural theorems, and no separate clean open statement is promoted. Those theorems are not reproduced in the supplied record.

The coefficient field, prime, central subgroup and exact notation $\mathrm{PC}$ are absent. For interpreting a possible $P_C$ notation only, if $p$ is fixed, $C$ is a central elementary abelian $p$-subgroup and coefficients are $\F_p$, multiplication $m:C\times G\to G$ induces a coaction
\[
 m^*:H^*(G;\F_p)\longrightarrow H^*(C;\F_p)\otimes H^*(G;\F_p).
\]
Its primitive/coinvariant subalgebra in the usual central-coaction convention is
\[
 P_C H^*(G;\F_p)=\{x:m^*(x)=1\otimes x\}.
\]
This is an explicitly conditional explanation, not a recovered definition of the excerpt's symbol. It differs from the Hopf-primitive identity $\Delta x=x\otimes1+1\otimes x$. No choice of maximal central subgroup, no restriction to $p$-groups and no additional quotient are inserted without the report's definition.

The dated literature triage checked 21 August 2026 says that the structural question is not resolved merely by bounds cited in the report, while the verification says structural answers immediately follow in the same abstract. These fields have different assessments and neither supplies a precise residual open assertion. The retained record therefore documents the extracted question and its limitations rather than inventing an unresolved structural conjecture or restating absent theorems.
\subsection{Short English statement}
This extraction turns a rhetorical question about primitive-central group cohomology into an open problem, but the supplied verification says the abstract immediately answers it. The structural theorems and exact coefficient/subgroup conventions are missing, so no separate clean conjecture is promoted.
\subsection{Sources}
\begin{itemize}
\item Local UnsolvedMath registry, permanent ID 30000320, exact source code catalogue entry; \emph{Rhetorical Primitive-Central Cohomology Question}. Proposed year 2005, registry status open, statement\_status unrecoverable. Supplied verification explicitly says structural theorems answer the rhetorical transition and retains it only for provenance.
\item Supplied report citation: \emph{Cohomology of Finite Groups: Interactions and Applications (2006)}, Oberwolfach Reports, DOI 10.4171/owr/2005/42, pp. 2375--2434, \url{https://doi.org/10.4171/owr/2005/42}. Title year differs from DOI/proposed year.
\item Supplied literature triage checked 2026-08-21, citing only ``original Oberwolfach report'', \url{https://doi.org/10.4171/owr/2005/42}. Its statement that bounds do not settle the structural question is retained alongside the verification's answered-rhetorical classification. The dated OWR-derived curation requests independent verification of time-sensitive status and cited claims; no separate precise residual conjecture is supplied.
\end{itemize}
\end{document}
